\chapter*{Удиртгал}

Компьютерын ухааны салбар үүсэн бий болохтой зэрэгцэн тооцооллын хүндрэл гэх нэр томъёо бий болсон ч
түүнээс ч өмнө үүний утга агуулга оршин байсан юм. Алан Тьюринг өөрийн алдарт өгүүллэгээ туурвиснаас
өдгөө 100 хүрэхгүй жилийн хугацаанд хүн төрөлхтөн асар хүчтэй тооцоолох машиныг бүтээсэн ба үүнийг асар
хурдтайгаар сайжруулсаар байгаа юм. Гэвч эдгээр хүчирхэг тооцоолох машины шийдвэрлэж чадахгүй олон асуудал байгаа
ба эдгээрийг одоогийн нөхцөлд шийдвэрлэж чадах сайн арга техник буюу алгоритм шаардлагатай байна. Иймд энэхүү
сэдэв нь компьютерын ухааны салбарын хамгийн чухал сэдвүүдийн нэг юм.

\section*{Судалгааны үндэслэл}

Компьютерын хүчин чадал асар хурдтайгаар хөгжиж байгаа ч бидний шийдвэрлэх шаардлагатай асуудлын хүндрэл түүнээс ч хурдтайгаар
нэмэгдэж байна. Иймд компьютер хэдий хурдтай хөгжиж байгаа ч энэ нь хэзээ ч хангалттай байж чадахгүй юм. Иймд тэдгээр асуудлыг
шийдвэрлэх үр ашигт алгоритм тооцооллын хүндрэлийн хамгийн том сэдвүүдийн нэг байсаар байна.

\section*{Судалгааны зорилго}

Энэхүү судалгааны ажлаар NP-hard ангийн "Явуулын Худалдаачны Асуудал"-ыг шийдвэрлэх өвөрмөц арга бүхий 3 төрлийн
алгоритмыг авч үзэн тэдгээрийг тооцооллын үр ашиг ба хэрэглээний хувьд харьцуулан судлах болно.

